% =========================================================================================
% Chapter 8 -- Discussion
%
% Per-module results live in Chapter 5, where each module's investigation ends in its own
% Results section. This chapter does NOT repeat them. It asks what they mean together:
% what the system can and cannot support, what failed and why, and what the numbers imply
% clinically. Cross-reference Chapter 5 rather than restating figures.
% =========================================================================================
\section{Evaluation and Discussion}
\label{sec:results}

\subsection{Introduction}

Each module was developed, measured and decided in Chapter~\ref{sec:perception}, where its
investigation ends in its own results section. This chapter does not repeat those figures. It
asks what the four results mean when the modules are assembled into one system: which claims
they support, which they do not, and what the assembled behaviour shows that no individual
metric can.

The distinction matters because the properties this project set out to demonstrate are not
properties of any single model. Whether an absent test is treated as absent, whether a
module's declared limits survive into the decision that follows, and whether the system
degrades safely when a component fails are all questions about the assembly.

\subsection{The Modules Side by Side}
\label{subsec:results-comparison}

Table~\ref{tab:modules-side-by-side} draws the four modules together. The headline metrics are
not comparable with one another --- a Dice coefficient, an AUROC, a balanced accuracy and a
macro F1 measure different things on different tasks --- and no conclusion should be drawn from
one number being larger than another. What can be compared is the distance of each result from
the baseline appropriate to its own task.

\begin{table}[H]
  \centering
  \small
  \caption{The four perception modules. Headline metrics are not comparable across rows;
    the distance of each from its own baseline is.}
  \label{tab:modules-side-by-side}
  \begin{tabularx}{\textwidth}{@{}lXlll@{}}
    \toprule
    Module & Task & Headline & Baseline & Cases \\
    \midrule
    Triage, reference & 3-tier urgency, tabular & macro F1 0.666 & 0.333 (chance) &
      16\,180 enc. \\
    Triage, deployed & the same, intake features only & macro F1 0.567 & 0.333 (chance) &
      16\,180 enc. \\
    Lung & 4 findings, multi-label & macro AUROC 0.788 & 0.500 (chance) & 165 \\
    Cardiac & segmentation $\rightarrow$ EF band & bal.\ acc.\ 63.4\% & 33.3\% (chance) &
      100 patients \\
    Gallbladder & 5 pathology classes & bal.\ acc.\ 59.7\% & 20.0\% (chance) & 199 \\
    \bottomrule
  \end{tabularx}
\end{table}

Read that way, the modules are more alike than their headline figures suggest. The gallbladder
classifier reaches roughly three times its chance level on a five-class problem; the cardiac
and triage models roughly twice theirs on three-class ones; the lung classifier sits well above
the uninformative 0.5 on every finding it models. None of them is close to the accuracy a
deployed clinical tool would require, and all of them are clearly doing something other than
guessing.

Triage appears twice because two models exist, and reporting only the higher figure would
describe a system that was never deployed. The reference model uses the registration fields a
hospital assigns; the deployed model uses only what a clinician types at the bedside, and it is
the one whose suggestions reach a screen. Section~\ref{subsec:triage-deployment} sets out the
full comparison. The single line worth carrying here is that the 0.099 of macro F1 separating
them is almost entirely recall on \textsc{low}: high-urgency recall is preserved, and the
dangerous miss --- a genuinely high-acuity patient graded low --- is marginally lower in the
deployed model than in the reference one.

The cardiac row needs a note, because the module produces two things. Its segmentation is
strong and close to the published benchmark --- Dice 0.928, 0.837 and 0.875 for left ventricle,
myocardium and left atrium against nnU-Net's 0.93, 0.86 and 0.89 --- while the clinical finding
derived from it, an ejection-fraction band, is the weaker figure quoted in the table. The
segmentation is not the agent's output; the band is, and it is the band a clinician would act
on. Reporting the stronger number as the module's result would overstate what the agent
delivers.

The same caution applies to plain accuracy, and it is the reason balanced accuracy appears in
the table for three of the four rows. On the cardiac cohort, 66 of 100 patients fall in one EF
band, so predicting that band for everyone scores 66\% while discriminating nothing --- higher
than the model's 65\%. Balanced accuracy of 63.4\% against a 33.3\% chance level is the honest
statement, and the same reasoning governs the triage and gallbladder figures.

Two observations survive the comparison. The first is that case counts are small everywhere
except triage, and Section~\ref{subsec:results-limitations} states what that forbids. The
second is that the triage model is the only one evaluated across multiple independent sources,
and its accuracy varies substantially between them --- from 0.617 to 0.789 across the three
datasets. A single pooled figure conceals that spread, which is why the per-source breakdown
is reported alongside it.

\subsection{Calibration as a System Property}
\label{subsec:results-calibration}

Four modules were calibrated by four different methods: isotonic regression for cardiac, Platt
scaling per finding for lung, temperature scaling for gallbladder, and Platt scaling for
triage. The choice was made per module rather than uniformly, because the shape of the
miscalibration differed. Temperature scaling has a single parameter and preserves the ranking
of classes, which suits a multi-class head that is confidently wrong at the top end. Isotonic
regression is non-parametric and can correct an arbitrary monotone distortion, which suits a
confidence derived from agreement across test-time augmentations rather than from a softmax.
Platt scaling fitted separately per finding suits a multi-label head whose four outputs are
miscalibrated to different degrees.

The system-level argument matters more than any of these choices. Calibration is what makes
four heterogeneous models comparable to one another. The reasoning layer receives a confidence
from a segmentation network, a confidence from a multi-label classifier and a confidence from
a tabular model, and it must weigh them against each other; without calibration those numbers
would sit on incommensurable scales and the escalation thresholds of
Section~\ref{subsec:reasoning-escalation} would have no defensible basis. A threshold of 0.80
means nothing unless 0.80 means approximately the same thing whichever module produced it.

Calibration is therefore not a finishing touch in this architecture. It is the precondition
for the architecture working at all, and the schema enforces it: a module that ships without
its calibrator reports \texttt{confidence\_calibrated: false}, and the reasoning layer
discounts it rather than treating an uncalibrated number as if it carried the same meaning.

The triage model illustrates both the value and the limit. Its expected calibration error fell
from 0.044 to 0.036 under calibration --- a real improvement, and a small one, because the
model was already reasonably calibrated. Calibration corrects miscalibration; it does not
manufacture reliability that was never there.

\subsection{What the Explainability Results Do and Do Not Establish}
\label{subsec:results-explainability}

The Grad-CAM statistics of Section~\ref{subsec:explainability} establish that attention is
concentrated rather than diffuse. That is evidence the models respond to localised image
content rather than to global brightness, framing or position, which is a real and
non-trivial thing to have shown: a classifier that had learned the average intensity of an
image would not produce concentrated attention maps.

It is not evidence that the attention falls on the anatomically correct structure. Establishing
that would require expert annotation of where the relevant structure lies in each image, which
this project does not have. The distinction is not pedantic. A model attending consistently to
a burned-in annotation marker, a probe artefact or a caption would also produce concentrated
maps, and the gallbladder module in particular showed a pronounced central prior that is
consistent with either a genuine anatomical focus or an acquisition convention.

Stating that boundary is worth more than the statistics themselves. Grad-CAM here is a
plausibility check that the models are not obviously broken, and it is reported as such rather
than as evidence of clinical validity.

\subsection{Clinical Reasoning Evaluation}
\label{subsec:results-reasoning}

The reasoning agent was evaluated on five synthetic scenarios, each designed to exercise a
specific failure mode: a case with key laboratory values absent, a case where triage and
imaging disagree, a case where the evidence agrees and the record is complete, a case where
every finding is negative, and a case where a requested organ was never scanned. Five cases is
a small benchmark, and that is stated here rather than deferred: no proportion computed over
five cases separates a real effect from a small one, and none is presented as if it did.

What was checked on each case was schema validity, whether the hallucination guard fired when
it should, whether every cited item existed in the state, and whether the escalation decision
matched the one the policy of Section~\ref{subsec:reasoning-escalation} specifies.

The result worth reporting is not a score but a change. With the evidence fields written as
free text, the model produced six unsupported-output errors across the five cases and the
safety layer withheld the differential on three of them. The failures were of a consistent
kind: an invented observation (``stable vitals'' for a patient whose vitals were not in the
state), a value relabelled against the record (a troponin the state printed as normal,
described as elevated), and --- once retrieval was added --- sentences drawn from the corpus
and placed among the patient's findings.

Constraining the interface removed them. When \texttt{supporting} and \texttt{contradicting}
were changed to accept only identifiers drawn from an enumerated list of citable facts, the
errors fell from six to zero and the withheld differentials from three of five to none. The
mechanism is worth stating precisely: the failures did not become easier to detect, they
became impossible to express. There is no identifier for a fact the state does not hold, none
for a test that was never performed, and none for a sentence out of a guideline.

Two faults resisted this treatment, both of them concerning how the answer is written rather
than what it claims: an abnormal value the model failed to cite, and a recommendation phrased
as though a test would settle a question rather than inform it. Two attempts to correct them
through the model failed --- first as an instruction in the prompt, then as a revision request
naming the exact identifier to add, across two correction rounds. They were resolved instead
by moving them out of the model: a comma-joined list is normalised in Python, and every
abnormal value is surfaced to the clinician from the state itself regardless of what the model
cited. The overclaiming wording was deliberately not corrected in code, because rewording a
recommendation changes what the system asserts, and it is reported as a warning instead.

The evaluation demonstrated the robustness of the safety and evidence-grounding mechanisms on
a five-case synthetic benchmark. It does not constitute clinical validation or a diagnostic
accuracy assessment, and no accuracy figure for the reasoning agent is reported anywhere in
this work, because no clinical ground truth exists against which one could be computed.

On retrieval, the finding is negative and is reported as such. On the five-case benchmark,
retrieval did not demonstrate a measurable improvement over the no-retrieval condition, and
one additional error occurred in the retrieval condition; its contribution to reasoning
quality remains inconclusive. Retrieval retains an architectural role --- it grounds the
model in a curated corpus whose provenance is recorded per passage --- but this project did
not measure a benefit from it.

\subsection{Overall System Evaluation}
\label{subsec:results-system}

Three properties appear only once the modules are assembled, and none of them is visible in
any per-module metric.

\paragraph{Graceful degradation.} A module that has no weights, no calibrator or no
implementation reports its status in the ordinary schema rather than raising or being absent.
The reasoning layer never branches on whether a key exists; an unsupported organ is an empty
finding list with a status. The consequence was tested directly with a deliberately failing
model backend: the escalation decision and its triggers were produced unchanged and the
differential was withheld with the backend error recorded. The safety-critical output survives
the failure of the least predictable component, which is the property the architecture was
arranged to have.

\paragraph{Propagation of declared limits.} Each module declares what it cannot exclude --- the
gallbladder classifier has no healthy class, the lung classifier does not model pneumothorax
--- and those declarations reach the escalation policy rather than stopping at the module
boundary. A scan in which nothing fires therefore escalates, on the ground that the models
cannot exclude disease, and does not read as a normal study. This is the inverse of the
behaviour a naive assembly would produce, in which an all-negative scan is the most reassuring
possible result.

\paragraph{Absence represented explicitly.} The system distinguishes three states that a
simpler design collapses into two: a finding detected, a finding screened for and not seen,
and a finding never assessed. The distinction is carried from the imaging module through the
clinical state into the prompt and the validators. Its practical effect was measured: a
defect in which an all-negative lung scan was represented by a placeholder entry in the
findings list caused the absence of findings to escalate less readily than their presence,
which is backwards, and the correction was verified end to end.

The behaviour on the worked scenarios of Section~\ref{subsec:integration-scenarios} exercises
all three together.

\subsection{Rejected Experiments}
\label{subsec:results-rejected}

Each rejection is reported inside the module that prompted it in
Chapter~\ref{sec:perception}. Table~\ref{tab:rejected} collects them so that the pattern is
visible.

\begin{table}[H]
  \centering
  \small
  \caption{Well-motivated changes tested and rejected on measurement.}
  \label{tab:rejected}
  \begin{tabularx}{\textwidth}{@{}lXl@{}}
    \toprule
    Change & Hypothesis and measured outcome & Decision \\
    \midrule
    External-dataset integration & Adding the external cardiac collection to training would
      bridge the generalisation gap; both integrations degraded in-domain performance without
      closing it & rejected \\
    Augmentation upgrade & Stronger augmentation would improve generalisation; the measured
      difference was within run-to-run variation & rejected \\
    USCL initialisation & Ultrasound-specific pretraining would beat ImageNet weights; it
      underperformed them & rejected \\
    Cardiac fine-tuning & Fine-tuning would close the external-domain gap; it improved the
      external set at the cost of the primary one & rejected \\
    Gallbladder taxonomy & An alternative grouping would separate the classes better; the
      collapsed confusion matrix predicted worse separation and training confirmed it &
      rejected \\
    Triage class weighting & Aggressive weighting would improve safety on the
      under-represented tier; it cost substantial accuracy for negligible safety gain &
      rejected \\
    Reason-for-visit grouping & A coarser grouping would generalise across sources; it lost
      accuracy, the fine-grained codes carrying signal the grouping discards & rejected \\
    \bottomrule
  \end{tabularx}
\end{table}

The table makes a point that no individual entry makes. Seven well-motivated changes were
tested and rejected on measurement rather than on preference, and each is reported with the
number that decided it. This is the methodological rule of Section~\ref{subsec:methodology}
shown working, and it is the strongest available evidence that the retained results were not
selected for being favourable. A project that reports only what worked gives a reader no way
to judge how much was tried.

The same discipline applied to the reasoning agent, where two changes were adopted and one
rejected: enumerated evidence was retained on measurement, a second correction round was
tested and reverted after it produced no change for a third more runtime, and the sharpened
complaint wording was retained only because it cost nothing, not because it was shown to work.

\subsection{Validating the Promotion Gate}
\label{subsec:results-mlops}

The gate described in Section~\ref{subsec:mlops} was exercised on real retraining runs rather
than on constructed metrics. The experiment isolates the variable the mechanism exists to
handle: two runs differ only in how much training data they were given, and the evaluation
folds are identical, so the comparison is like for like. Training volume is reduced by
dropping training \emph{case groups} rather than frames, because removing frames would leave
the same patients with thinner coverage, which is not what a smaller cohort means clinically.

\begin{table}[H]
  \centering
  \small
  \caption{Identical evaluation folds; only the quantity of training data differs. The
  candidate improves on aggregate performance and on two of three guarded findings.}
  \label{tab:mlops-ab}
  \begin{tabular}{@{}lrrrl@{}}
    \toprule
    Metric & Half data & Full data & Change & Gate \\
    \midrule
    Macro-F1                  & 0.4303 & 0.5004 & $+0.0701$ & pass \\
    Recall, B-lines           & 0.7302 & 0.8413 & $+0.1111$ & pass \\
    Recall, consolidation     & 0.7167 & 0.7500 & $+0.0333$ & pass \\
    Recall, pleural effusion  & 0.6250 & 0.5833 & $-0.0417$ & \textbf{fail} \\
    \midrule
    \multicolumn{4}{@{}l}{Decision} & \textbf{rejected} \\
    \bottomrule
  \end{tabular}
\end{table}

Doubling the training data raised macro-F1 by seven points and recall on the commonest finding
by eleven, and lowered recall on the rarest one. A gate that compared headline performance
would have accepted this candidate and deployed a model that misses more pleural effusions
than the one it replaced. This is the failure mode the guarded metrics exist for, and it
arose from a routine change in training data rather than from an adversarial example.

The correct statement of the result is that the full-data model is better on aggregate
performance and \emph{not eligible for promotion under the stated criteria} --- not that it is
a better model. The distinction is the entire contribution of the mechanism.

Four further candidates were rejected during the same work: one whose evaluation protocol
differed from the incumbent's and was therefore not comparable, one that fell below the
macro-F1 floor, and two whose headline metric regressed. No candidate was accepted. The
guard tolerance was not relaxed to produce one, because a threshold revised after seeing the
result it decides is a description of that result rather than a criterion.

\subsubsection{A limitation the experiment exposes}

Pleural effusion has 24 positive clips in the evaluation cohort, so a single clip changing
class is approximately 4.2 percentage points of recall --- and the observed $-0.0417$ is one
clip. The guard tolerance of two points is therefore finer than the data can resolve, and the
gate cannot distinguish a real regression on this finding from sampling noise.

This argues for making the tolerance sensitive to class support, or for a larger case-grouped
evaluation cohort, and it argues against relaxing the threshold on the evidence of the run
that failed it. It is recorded here rather than acted on for that reason.

\subsubsection{Optimistic bias in the original lung figure}

Extracting the training pipeline from the notebook surfaced a methodological problem in the
original evaluation. The notebook applies early stopping and retains the best checkpoint by
AUROC measured on the cross-validation fold that it then reports as held-out performance.
Model selection and performance reporting therefore use the same data, which biases the
reported figure upwards by an amount that cannot be recovered after the fact.

This is not label leakage in the formal sense --- no training example carries a test label ---
and it should be described precisely as \emph{optimistic bias from checkpoint selection on the
reported evaluation fold}. Its practical effect is that the macro-F1 of 0.5475 quoted for the
lung model in Section~\ref{subsec:lung-results} is not directly comparable with figures
produced by a protocol that holds the scored fold aside.

The extracted pipeline takes its stopping signal from an inner validation split held out by
case group from the training folds, so the scored fold is not consulted until it is scored.
This is stricter, produces lower numbers, and is the protocol under which all subsequent
candidates are evaluated. The two protocols are recorded as distinct, and the gate's
same-evaluation-set check refuses to compare across them --- which is how the difference was
detected in the first place.

\subsection{Strengths}
\label{subsec:results-strengths}

Four claims are defensible from measurement rather than from intent, and each is checkable by
an examiner.

\paragraph{Case-level evaluation throughout.} Three of the four datasets provide image-level
or frame-level labels only. Splitting on those directly leaks the same patient across train
and test and inflates every figure. Case-level grouping was reconstructed for each --- for the
lung dataset by a filename heuristic whose residual uncertainty is stated in
Section~\ref{subsec:results-limitations} --- and every reported metric is computed over
independent cases.

\paragraph{Calibrated confidence carried end to end.} Confidence is calibrated in each module,
exported with the weights, and consumed by the reasoning layer, rather than being computed and
discarded at the module boundary. Where a calibrator is missing the output says so.

\paragraph{Absent evidence represented explicitly.} The three-state distinction described in
Section~\ref{subsec:results-system} is enforced by the schema, not by convention: a report
claiming an assessment occurred while recording nothing about it is rejected.

\paragraph{Safety decisions independent of the language model.} Escalation is computed from
the structured state before the model is consulted and is unchanged by anything the model
does. This is demonstrated rather than asserted: 339 tests exercise the safety layer with no
model present, and the decision was verified identical across three different backends
including one that fails.

\subsection{Limitations}
\label{subsec:results-limitations}

Each limitation below is stated as a bound on a specific claim. A limitation that does not
change what may be concluded is not listed.

\paragraph{Small case-level datasets.} 165 lung cases, 199 gallbladder cases, 100 cardiac
validation patients. Confidence intervals on every per-module figure are wide, and differences
of a few percentage points between configurations are not interpretable. This forbids
comparative claims about closely-spaced variants.

\paragraph{Public collections rather than consecutive clinical series.} The images are
teaching-atlas stills and curated public clips, which are better framed and less noisy than
bedside acquisition. This forbids any claim that the reported accuracy would transfer to
point-of-care use.

\paragraph{No external validation for most modules.} Only the triage model was evaluated
across independent sources, and its accuracy varied from 0.617 to 0.789 between them. For the
imaging modules the absence of external evaluation forbids any claim of generalisation beyond
the collection each was trained on.

\paragraph{The lung case mapping is a heuristic.} Case identity was reconstructed from
filenames rather than provided by the dataset. Residual uncertainty remains, and it forbids
treating the lung case count as exact.

\paragraph{The cardiac module is domain-limited.} It was trained on CAMUS-like four-chamber
stills and its performance dropped on an external collection; fine-tuning to close that gap
was tested and rejected because it cost primary-domain accuracy. This forbids applying it to
acquisitions unlike its training domain.

\paragraph{No clinical outcome validation and no prospective evaluation.} Nothing in this
work was measured against what happened to a patient, or against what a clinician concluded.
This forbids every claim of clinical utility, and is the reason no diagnostic accuracy figure
appears for the system as a whole.

\paragraph{The reasoning benchmark is five synthetic cases.} It measures whether the safety
mechanisms behave as specified. It forbids any statement about how the reasoning agent would
perform on real presentations, and it is too small to establish that retrieval does or does
not help.

\paragraph{The per-organ inference functions were not under test, and now are.} This limitation
was previously total: every test built module reports directly, none imported the per-organ
prediction code, and the suite therefore proved the contract was \emph{enforced} without proving
that any module \emph{honoured} it.

It has been closed for all three organs, whose prediction code was extracted from the notebooks
and is now exercised by sixteen tests running real inference on the shipped checkpoints
(Section~\ref{subsubsec:agent-extraction}).

The reason to report this as a result rather than as housekeeping is what closing it found. Four
defects surfaced within the extraction, and not one of them raised an error: a decision threshold
that silenced the lung module so that every study returned negative; an ejection fraction
computed from a single frame; a failed systolic segmentation entering the formula as an area of
zero and emerging as a normal ventricle at maximum confidence; and a missing image raising an
exception instead of degrading to a failed report.

Three of the four produced a \emph{reassuring} output from a broken measurement. That is the
direction that matters in a screening system, and it is invisible to any amount of inspection ---
each looked like a working module returning a plausible number. They became visible when the code
became importable and a test asserted a property of it. The generalisable claim is narrow and
worth stating: research code that cannot be imported cannot be tested, and the failures it hides
are selected for being quiet.

\subsection{Clinical Interpretation}
\label{subsec:results-discussion}

None of the four modules is accurate enough to be relied upon alone, and the gallbladder
classifier at 59.7\% balanced accuracy on five classes is the clearest example: as a
standalone diagnostic it would be unusable. Read as one input among several, weighted by a
calibrated confidence, and unable to assert that a patient is healthy because it has no
healthy class, it is a different proposition --- not because the number is better than it
looks, but because the architecture never asks it to carry a decision by itself.

That is the argument this project is really making. A well-calibrated model of moderate
accuracy can be more useful inside a system like this than a more accurate model whose
confidence cannot be trusted, because the reasoning layer can act on the first and cannot
safely act on the second. A confidence of 0.4 that means 0.4 tells the escalation policy to
escalate; a confidence of 0.9 that means 0.6 tells it to proceed, and it is wrong. The
project's emphasis on calibration follows from that, and so does its willingness to report a
model that is only moderately accurate rather than tuning until a headline figure improves.

Returning to the four difficulties of Section~\ref{subsec:problem}:

\paragraph{Operator dependence and image quality} is \emph{not addressed}. Every model here
was trained and evaluated on curated images. The difficulty is acknowledged in the corpus,
which records that a point-of-care study is dependent on the skill of the clinician
performing it, but nothing in this work measures or mitigates it.

\paragraph{Fragmented evidence across modalities} is \emph{addressed}. The clinical state
assembles imaging, vitals, laboratory results and triage into one structure with a single
contract, and the reasoning layer reads all of it together. This is the difficulty the
architecture most directly answers.

\paragraph{Confidence that cannot be compared across models} is \emph{addressed}, by
calibration carried end to end and by the schema's refusal to let an uncalibrated number pass
as a calibrated one.

\paragraph{Unreliable reasoning over the assembled evidence} is \emph{mitigated, not solved}.
Unsupported evidence citations were eliminated by constraining the interface, and the safety
decision does not depend on the model at all. But the model still reasons from part of the
available evidence rather than all of it, still phrases recommendations more strongly than
decision support should, and has not been evaluated against clinical ground truth. What this
work establishes is that a language model of this size can be constrained so that its failures
are contained and visible --- not that its reasoning is reliable.
